\chapter{Implementation and Experimental Setup}
\label{ch:implementation}

\Cref{ch:benchmarking-methodology} fixed the protocol. This chapter describes the code that
applies it, as the repository \texttt{fil2405/nesy-nids-bench} holds it. Every number comes
from the data files, the configurations or the result files. The plan of the thesis lists
components that the repository does not contain: an LSTM, a semantic-loss variant, window
features and a temporal split beyond the two days of CIC-DDoS2019. They are not described as if
they existed, and \cref{sec:impl-summary} lists them. The real-time pipeline belongs to the
co-supervisor and appears only as far as it is fixed.

% ================================================================
% 9.1 Hardware and software environment
% ================================================================
\section{Hardware and software environment}
\label{sec:impl-environment}

The code is a Python package, \texttt{nesybench}. The networks are written in TensorFlow
\citep{abadi2016tensorflow} through the \texttt{tf\_keras} package, because version 2.1 of
\texttt{ltn}, the library for Logic Tensor Networks \citep{badreddine2022ltn}, needs Keras 2.
Random Forest comes from scikit-learn \citep{pedregosa2011scikit} and boosting from XGBoost
\citep{chen2016xgboost}. Tables are pandas frames, and the cleaned CIC datasets are cached as
Parquet files.

Two environments were used (\cref{tab:impl-environments}). The server ran every run that
trains a model: the searches of hyperparameters, the final runs and the stress runs. The laptop
ran the development, the evaluations of saved models (E10, E12, E17 and the post-check) and the
timings of E13.
% FIXME(open): the NSL-KDD Random Forest and XGBoost use the configuration of the first search (tuning_round1), which ran on the laptop; the other arms were searched on the server. Confirm, and if so state the exception here and in Ch. 8, Sec. 8.1.5.
% FIXME(open): saved_runs.sh launches latency.py on the CPU, but saved_eval.py and e13_time.sh use the default device and do not record it, and the laptop has a GPU that TensorFlow sees. Confirm that the evaluations of saved models and the E13 timings ran on the CPU, as Ch. 8, Sec. 8.1.5 states, or correct Ch. 8.

\begin{table}[t]
\centering
\footnotesize
\caption{The two environments. The versions of the server are those recorded in its result files.}
\label{tab:impl-environments}
\setlength{\tabcolsep}{4pt}
\begin{tabular}{@{}p{2.4cm}p{5.2cm}p{5.4cm}@{}}
\toprule
 & Server & Laptop \\
\midrule
Used for & Tuning, final and stress runs & Development; evaluations of saved models; E13 timings \\
System & Linux, kernel 7.0.0-30 & Debian under WSL2 on Windows 11 \\
Hardware & not recorded & Core i7-12650H, RTX 3050 Laptop with 4 GB, 12 GB of RAM for WSL2 \\
Python & 3.14.4 & 3.11.2 \\
TensorFlow & 2.22.0-rc0 & 2.21.0 \\
NumPy, scikit-learn & 2.5.4, 1.9.1 & 2.4.6, 1.9.1 \\
XGBoost, \texttt{ltn} & 3.4.1, 2.1 & 3.2.0, 2.1 \\
Processes & seeds and trials in parallel & one at a time \\
\bottomrule
\end{tabular}
\end{table}
% FIXME(open): the hardware of the server (processor, memory, GPU) is in no result file; ask the co-supervisor and fill the cell.

The neural arms train on the GPU and Random Forest on all cores of the processor. On the server
the seeds of a final or stress run and the trials of a search run as parallel processes that
share the GPU. These are independent runs: no model is trained across processes or nodes.

A run seeds Python, NumPy and TensorFlow and switches on the deterministic operations of
TensorFlow before anything is built. In our checks, repeating a run of the MLP, of the two LTN
arms or of XGBoost with the same seed in one environment gave the same validation history and
the same selected epoch. The Random Forest was not checked. The deterministic operations slow the
GPU down, and every training time of this thesis was taken with them on. The server runs a
TensorFlow release candidate, so the searches, the final runs and the stress runs rest on a
pre-release. Across environments and versions we claim nothing, not even the last digit of a
probability.

The two environments are kept apart by three rules. The cells of one table come from one
environment, and each table of \cref{ch:results} names it. Times are compared within one
environment only. An evaluation on the laptop of a model trained on the server prints the
macro-F1 stored by the run beside the one it recomputes, so a difference would show.
The environment of the real-time pipeline is the co-supervisor's and is not described here
(\cref{sec:realtime-pipeline}).
% FIXME(open): hardware and software of the machine that runs the streaming pipeline.

% ================================================================
% 9.2 Data preparation
% ================================================================
\section{Data preparation}
\label{sec:impl-data}

Each dataset has a loader that returns the raw table with the fine label and the class. A
scaler and a one-hot coding are fitted on the training part only. The raw rows travel with the
transformed ones, so the rules read the units of the data.

\subsection{NSL-KDD pipeline (inherited, corrected)}
\label{sec:impl-nslkdd}

NSL-KDD is the version of KDD Cup 99 that Tavallaee et al. \citep{tavallaee2009kddcup99} made
by removing redundant records. We use the official files, KDDTrain+ with 125\,973 records and
KDDTest+ with 22\,544. A record has 41 features, an attack name and a difficulty score, which
comes from classifiers that had seen the labels and is never an input.

The mapping of attack names to five classes, the one-hot coding of \texttt{protocol\_type},
\texttt{service} and \texttt{flag} (3, 70 and 11 values) and the standardisation of the 38
numeric columns come from the earlier project of the author. The network has 122 inputs. Three
things differ. The validation part is cut from the training file, so the test file chooses
nothing. No oversampling is applied. The two \texttt{worm} records of KDDTest+ are kept, in DoS.

Three attack names occur only in the test file and are assigned by convention:
\texttt{snmpgetattack} and \texttt{httptunnel} to R2L and \texttt{worm} to DoS. With
\texttt{httptunnel} in U2R, KDDTest+ would have 200 U2R and 2\,752 R2L records, not 67 and
2\,885. In KDDTest+, 3\,750 records, 16.6\,\%, have an attack name that KDDTrain+ lacks, and
R2L changes its composition: 890 of its 995 training records are \texttt{warezclient}, which the
test file does not contain.
% FIXME(open): sources for the family of snmpgetattack, worm and httptunnel are secondary (arXiv 1910.01114, 1803.06762, 2112.02080); add them to the bib if the sentence is to carry a citation, and close the counts FIXME of Ch. 3, Sec. 3.4.3 with this section.

\subsection{CIC-IDS2017 / CIC-DDoS2019 pipeline and the corrected-label variant}
\label{sec:impl-cic}

\paragraph{CIC-IDS2017.}
The original capture \citep{sharafaldin2018cicids2017} has known errors in flow extraction and
labelling \citep{engelen2021troubleshooting}. Liu et al. \citep{liu2022error} publish a later
version with improved ground truth and more features, and we use it. It ships five daily files
with 2\,099\,976 flows: identifiers, addresses, ports, the protocol, a time stamp, the numeric
flow features, a label and an Attempted category. An attack flow that carried no data forward is
labelled ``Attempted''. The authors recommend counting these as benign, and we do so. After
cleaning they are 9\,106 of the Normal flows.
% FIXME(open): the recommendation on the Attempted flows is on the download page of the improved version, not in the abstract read of the paper; read the paper and cite the section, or reword to the page.

The cleaning is done once, before any split. Five flows with infinite rates go, all Normal. So
do 384\,630 exact duplicates on features and label, 223\,023 of them Probe and 161\,607
Normal, which leaves 1\,715\,341 flows; the launcher of the final runs checks this count before
it starts. Probe falls from 230\,833 flows to 7\,810, since 97\,\% of the port-scan flows differ
only in the port scanned, which the network does not see. The 256 rows with equal features and
two labels are kept.

\begin{table}[t]
\centering
\footnotesize
\caption{Normal flows per day, as read from the files and after cleaning, and the attack flows
that remain. A day is a day of a laboratory network.}
\label{tab:impl-benign-day}
\setlength{\tabcolsep}{4pt}
\begin{tabular}{@{}lrrp{5.6cm}@{}}
\toprule
Day & Normal, as read & Normal, cleaned & Attack flows after cleaning \\
\midrule
Monday & 371\,624 & 350\,820 & none \\
Tuesday & 315\,145 & 290\,135 & R2L 6\,933 \\
Wednesday & 324\,996 & 288\,356 & DoS 171\,645 \\
Thursday & 290\,169 & 254\,582 & Probe 6\,127, R2L 104, U2R 36 \\
Friday & 292\,611 & 249\,040 & DDoS 95\,144, Probe 1\,683, Botnet 736 \\
\midrule
CIC-DDoS2019, 12 Jan & 56\,863 & --- & DDoS \\
CIC-DDoS2019, 11 Mar & 56\,965 & --- & DDoS \\
\bottomrule
\end{tabular}
\end{table}

\Cref{tab:impl-benign-day} gives the benign flows of each day for the alert budget per day.
Which of the two counts is the day of the budget is not fixed.
% FIXME(open): the alert budget per day of Ch. 8 uses the benign flows of a day; decide whether it is the count as read or the cleaned one and say so in Ch. 8, Sec. 8.2.3.
% FIXME(open): the figures of the two CIC-DDoS2019 days are the benign rows of the CSV files, not the 45143 benign rows of the sampled test day.

\paragraph{Mechanism axioms on single flows.}
The labelling of this version also decides what an axiom built on an attack mechanism can do on
a single flow. Three such axioms were written with thresholds fixed before the data were read,
with a precision of 0.90 as the criterion of acceptance, and measured on the training and
validation parts of seed 0, 1\,372\,272 flows. A slow DoS axiom reached a precision of 0.008 and
covered 4\,\% of the Slowloris flows. A scan axiom, a TCP flow with at most two forward packets,
at most one backward packet and no forward payload, fired on 9\,720 flows with a precision of
0.092 and covered 14\,\% of Probe; 7\,771 of its flows were BENIGN and 829 were Botnet flows
labelled Attempted. A half-open axiom, a TCP flow with a SYN flag and no backward packet, fired on
5\,233 flows with a precision of 0.000. It hit no DoS flow, but 83\,\% of the DoS Slowhttptest
flows labelled Attempted and 89\,\% of the DoS Slowloris ones. In our reading, the axiom
describes the attempt, while the labels mark as Normal every attempt that had no effect. The
knowledge base of CIC-IDS2017 therefore holds
no mechanism axiom, and its two axioms are extracted from the data (\cref{app:axioms}).

\paragraph{CIC-DDoS2019.}
The dataset \citep{cic2019ddos} has a training day, 12 January, with 11 CSV files and
50\,063\,112 flows, and a test day, 11 March, with 7 files and 20\,364\,525. Of the
70\,427\,637 rows only 113\,828, 0.16\,\%, are benign. The name of a file is not its label:
\texttt{UDPLag.csv} of the test day holds 606\,749 Syn flows and 1\,873 UDPLag. The loader
reads the label from the column. We found no corrected variant of this dataset.

The days are too large to train on, so a cap on the labels fixes the sample. Before any
cleaning, the loader keeps every benign row and at most 90\,000 rows of each attack label in the
training day and 30\,000 in the test day, drawn at random with a fixed seed in two passes over
the archives. The label \texttt{UDP-lag} is unified with \texttt{UDPLag}, and the header length
that CICFlowMeter writes twice is dropped after a check that the copy is equal. Rows with
infinite values and duplicates on features and label are then dropped. The training day goes
from 1\,047\,302 sampled rows to 340\,560, after 36\,696 non-finite rows and 670\,046
duplicates, and 14.8\,\% of them are Normal. The test day goes from 238\,838 rows to 79\,083,
after 7\,708 and 152\,047, and 57.1\,\% are Normal. Rows with equal features and two labels
are kept: 22\,469 in the training day and 1\,191 in the test day.

The days share few labels. Only Syn and UDPLag occur in both. DNS, NTP, SNMP, SSDP and TFTP
are absent from the test day, and Portmap, with 588 flows, occurs only there, so a recall on it
is a rough estimate. LDAP, MSSQL, NetBIOS and UDP occur only in the test day under these names,
and in the training day with the prefix \texttt{DrDoS\_}. The names with and without the prefix
are not merged.

\subsection{Common feature schema for cross-dataset experiments}
\label{sec:impl-schema}

NSL-KDD has 41 connection features and shares no feature space with the CIC datasets, so it
takes no part in a transfer. The two CIC datasets come from CICFlowMeter, in two versions that
spell many columns differently. The code builds their common schema as a dataset of its own,
\texttt{cic\_transfer}, used by the transfer of E9.

A column name is normalised by lowering its case, dropping every character that is not a letter
or a digit and reducing the plurals ``packets'' and ``bytes'', and two columns match when their
normalised names are equal. Sixteen pairs that normalisation does not match are renamed by
hand. Among them are the count of backward packets, the minimum and maximum packet length, the
six bulk averages, the initial window bytes in both directions and the count of CWR flags,
which the header of CIC-DDoS2019 misspells as CWE. The common columns take the names of
CIC-DDoS2019. The loader stops if two columns of one side map to one name or if the protocol is
not among the common columns. The schema has 77 columns, the same in the training and the test
part. Identifiers, addresses, ports, time stamps, the Attempted category of CIC-IDS2017 and the
column \texttt{Inbound} of CIC-DDoS2019 stay out, as they do in each dataset alone.

Training and validation come from CIC-IDS2017: the rows of Normal, the Attempted flows included,
and those of DoS and DDoS, merged into one class Attack. The rows of Probe, R2L, U2R and Botnet
are dropped. The test part is the test day of CIC-DDoS2019, with DDoS as Attack. The task is
binary, with 1\,432\,933 Normal and 266\,789 Attack rows for training and validation, and
45\,143 and 33\,940 in the test part. Rows are not deduplicated again after the projection onto
the common columns. Two knowledge bases are written in the names of the target. The first holds
the two axioms of CIC-IDS2017 with target Attack. The second adds the UDP axiom of
CIC-DDoS2019, a UDP flow with no backward packet, whose threshold was chosen on the training day
of that dataset. Since the extractors differ, a loss on the target mixes a change of attacks
with a change in how the features are computed.

The class schema is that of \cref{tab:impl-classes}, in the same names and order wherever a
class exists. Outside the transfer the two CIC datasets share only Normal and DDoS.

\begin{table}[t]
\centering
\footnotesize
\caption{Classes in the order of the class index, and the fine labels of CIC-IDS2017 with
their flows after cleaning.}
\label{tab:impl-classes}
\setlength{\tabcolsep}{4pt}
\begin{tabular}{@{}lp{10.6cm}@{}}
\toprule
Class & Fine labels \\
\midrule
\multicolumn{2}{@{}l}{\emph{NSL-KDD:} Normal, DoS, Probe, R2L, U2R} \\
\multicolumn{2}{@{}l}{\emph{CIC-DDoS2019:} Normal, DDoS (every attack label, kept in a column)} \\
\multicolumn{2}{@{}l}{\emph{Transfer:} Normal, Attack (DoS and DDoS of CIC-IDS2017; DDoS of CIC-DDoS2019)} \\
\midrule
\multicolumn{2}{@{}l}{\emph{CIC-IDS2017}} \\
Normal & BENIGN 1\,423\,827; Attempted of every attack label 9\,106 \\
DoS & Hulk 158\,468; GoldenEye 7\,567; Slowloris 3\,859; Slowhttptest 1\,740; Heartbleed 11 \\
DDoS & DDoS 95\,144 \\
Probe & Portscan 1\,683; Infiltration\,--\,Portscan 6\,127 \\
R2L & FTP-Patator 3\,972; SSH-Patator 2\,961; Web Brute Force 73; XSS 18; SQL Injection 13 \\
U2R & Infiltration 36 \\
Botnet & Botnet 736 \\
\bottomrule
\end{tabular}
\end{table}

For CIC-IDS2017 the classes follow the families of \cref{ch:attacks-ddos}, with DDoS, many
sources, apart from DoS, one source. Infiltration as U2R and Heartbleed as DoS are our
conventions, and no source settles them. R2L collects the password guessing on FTP and SSH
with three web attacks, so it is not remote-to-local in the sense of the KDD families. The same
scan carries two labels, Infiltration\,--\,Portscan on Thursday and Portscan on Friday, and both
are Probe. Each attack class lives in one or two days (\cref{tab:impl-benign-day}), so a split
by day would remove whole classes from training. The split is random and stratified, and flows
of one attack session fall in training and in test, which limits the result
(\cref{sec:threats}).

\paragraph{Columns that are not inputs.}
In CIC-IDS2017 the identifiers, addresses, ports, time stamp, Attempted category and day are
not inputs. They stay in the raw table, where the rules can read them. They are capture
artefacts: in the cleaned file the DDoS flows come from one source address to one destination
address and one port, while the Normal flows come from 174 sources to 19\,044 destinations.
CIC-DDoS2019 is treated alike, and its column \texttt{Inbound}, a flag of the capture
direction, is excluded too: in the training day 99.8\,\% of the DDoS rows have it equal to 1
and 85.0\,\% of the Normal rows equal to 0. What its exclusion costs was not measured.

\paragraph{Protocol.}
The protocol is an input in all three datasets, one-hot, and the same for every arm. In
CIC-IDS2017 it is a shortcut: every Botnet, DDoS, DoS, R2L and U2R flow is TCP, as are 7\,629
of the 7\,810 Probe flows, while 844\,318 Normal flows, 58.9\,\%, are not. In CIC-DDoS2019 the
mix changes between the days: 261\,613 of the 290\,011 DDoS rows of the training day are UDP,
against 28\,211 of 33\,940 in the test day. A result that rests on the protocol is no evidence
about the knowledge.
% FIXME(open): keeping the protocol as an input is in the code but not yet confirmed as a decision. Confirm, and extend the threats of Ch. 8, Sec. 8.8.
% FIXME(open): the number of network inputs of CIC-IDS2017 and CIC-DDoS2019 is in no table; read it from the feature names of one run. The parameter count of the 1D-CNN on CIC-IDS2017 implies 84 or 85 inputs.

No window or source statistic is computed, as \cref{ch:usecases-design} states.

\subsection{Split strategy: train / validation / test, temporal and random variants}
\label{sec:impl-splits}

Two strategies exist, both stratified by class. The official-test strategy keeps the test
file and cuts a validation part of 15\,\% from the training file with the seed of the run. It
serves NSL-KDD, CIC-DDoS2019 and the transfer. The random strategy cuts one table in three. It
takes the test part first, 20\,\% of the rows, with a fixed test seed of 0, and then validation,
20\,\% of the rows, from the rest with the seed of the run. It serves CIC-IDS2017, so every arm
and seed meets the same test. Splits are not saved but rebuilt from the seed. Each run writes
the rows per class of its training, validation and test parts to its result file, in the
field \texttt{split\_sizes}. The final runs lack the field, and their sizes are rebuilt from
the seed.

A temporal split exists only as the two days of CIC-DDoS2019, and across datasets as the
transfer. NSL-KDD has no time stamps, and CIC-IDS2017 has no split by day. The code offers
instead four options on the parts that the strategy returns. A fraction of the training part is
kept, stratified. Chosen classes of the training part are thinned at random. A floor of rows
per class applies to both: a class never falls below it, or below what it had, and rows dropped
by the fraction are put back at random until the floor is met. With the option
\texttt{shrink\_val} the validation part receives the same cut, with the seed of the run plus
one. The test part is never cut.

E7 uses the fraction and E8 the thinning, both with a floor of five rows and the validation part
cut. On NSL-KDD with seed 0, the level of 5\,\% trains on 2\,862, 1\,952, 495, 42 and 5 rows of
Normal, DoS, Probe, R2L and U2R, 5\,356 in all, and validates on 505, 344, 87, 8 and 5. The
whole validation part would have 18\,896 rows. The search on CIC-IDS2017 thins Normal to a
quarter, with no floor and the validation part whole.

The random strategy also holds labels out for E9. The listed fine labels leave the table
before the split, and with the option \texttt{holdout\_attempted} so do their flows labelled
Attempted. The test part is cut from what remains, and the held-out rows are added to it, the
Attempted flows with their class, Normal. A label that would empty its class is refused. The
test part of E9 is thus cut without the held-out rows and is not the test part of the final
runs.

\begin{table}[t]
\centering
\footnotesize
\caption{Rows of each part for seed 0 at full size. The parts of the CIC datasets are derived
from the class totals with the stratified split of the code, and the NSL-KDD validation part is
the one recorded by the runs. Another seed moves rows between training and validation only.}
\label{tab:impl-splits}
\setlength{\tabcolsep}{5pt}
\begin{tabular}{@{}llrrr@{}}
\toprule
Dataset & Rows & Training & Validation & Test \\
\midrule
NSL-KDD & R2L & 846 & 149 & 2\,885 \\
 & U2R & 44 & 8 & 67 \\
 & all & 107\,077 & 18\,896 & 22\,544 \\
\midrule
CIC-IDS2017 & R2L & 4\,222 & 1\,407 & 1\,408 \\
 & Botnet & 442 & 147 & 147 \\
 & U2R & 22 & 7 & 7 \\
 & all & 1\,029\,204 & 343\,068 & 343\,069 \\
\midrule
CIC-DDoS2019 & Normal & 42\,967 & 7\,582 & 45\,143 \\
 & all & 289\,476 & 51\,084 & 79\,083 \\
\bottomrule
\end{tabular}
\end{table}

\Cref{tab:impl-splits} gives the sizes. U2R has 8 validation rows on NSL-KDD and 7 on
CIC-IDS2017, so one row moves its recall by an eighth or a seventh. The selection criterion of
\cref{ch:benchmarking-methodology} damps this and does not remove it. R2L is 0.8\,\% of the
NSL-KDD validation part and 12.8\,\% of the test file, and mostly \texttt{guess\_passwd} and
\texttt{warezmaster} there. The validation part carries no such shift.
% FIXME(open): the per-class sizes of the CIC parts in tab:impl-splits are derived from the class totals; the totals agree with the 1372272 training and validation flows of CIC-IDS2017. Check the rows of R2L, Botnet, U2R and Normal against split_sizes of seed 0 of an E4 run (NSL-KDD, CIC-IDS2017) and of a full-size run of CIC-DDoS2019.

% ================================================================
% 9.3 Models under test
% ================================================================
\section{Models under test}
\label{sec:impl-models}

A model is a class with three methods: \texttt{fit} on the splits and the configuration,
\texttt{predict\_proba} on a part, and \texttt{save}. The registry has seven names,
\texttt{rule\_only}, \texttt{rf}, \texttt{xgb}, \texttt{mlp\_ce}, \texttt{cnn1d},
\texttt{ltn\_sup} and \texttt{ltn\_full}, and \texttt{mlp\_ce} runs with or without class
weights, which makes eight arms. All arms get the same split, the same transformed matrix and
the same seed. The base configuration of each arm holds starting values. The search of
\cref{sec:impl-tuning} imposes the common network of \cref{ch:benchmarking-methodology} on the
four neural arms and its own network on the 1D-CNN, and writes the tuned configuration, which
the final runs use. Each neural run records its number of parameters. The LSTM and the
semantic-loss variant are not implemented, since the registry has no entry for them.

\subsection{Rule-only symbolic classifier (knowledge without learning)}
\label{sec:impl-rules}

The arm reads a knowledge base in YAML and learns nothing. An axiom has an identifier, a group,
an origin, a target class and a list of conditions on raw columns, each an equality with a
category or a comparison of a column with a threshold $\tau$. The antecedent is the product of the
conditions with hard thresholds, and the axiom fires when the product is at least 0.5. The arm
predicts the target of the axiom that fires. When several fire, the first in the order of the
file wins. When none fires, the flow goes to Normal. The output is a vector with a single one,
so the arm has no scores between 0 and 1. The fit only counts the training flows that the
axioms cover and those on which they conflict, and the saved arm is a copy of the knowledge
base.

A class without an axiom is never predicted, except Normal, which is the default. In the main
knowledge bases this holds for R2L and U2R on NSL-KDD, for Probe, R2L, U2R and Botnet on
CIC-IDS2017, and for every class but DDoS on CIC-DDoS2019. The variants of \cref{app:axioms}
add axioms for some of these classes.

\subsection{Classical ML: Random Forest, XGBoost}
\label{sec:impl-trees}

Random Forest is the scikit-learn classifier, with all cores and the seed of the run. It has no
epochs, so nothing is selected on the validation part. XGBoost uses the multiclass objective
with softmax probabilities and the histogram method, for at most 300 rounds. The validation part
drives the same stop as in the neural arms: the mean macro-F1 of the last five rounds, patience
ten, and the rounds after the best one are not used at prediction. Both trees take class weights
as sample weights, proportional to the count of the class raised to minus the power and scaled
so that the mean weight over the training rows is 1. The power is searched, and 0 means no
weights. The validation part is not weighted. Both trees read the same standardised and one-hot
matrix as the networks.

\subsection{DL baselines: MLP, 1D-CNN, [EXT] LSTM}
\label{sec:impl-mlp}

The baseline is a multilayer perceptron: dense layers of 256 and 128 units with ReLU, each
followed by dropout 0.2, and a softmax output. The weights start from Glorot uniform draws that
depend on the seed. Training uses Adam, batches of 256 and the cross-entropy, for at most 100
epochs. After each epoch the macro-F1 of the validation part is computed, and the selection
rule of \cref{ch:benchmarking-methodology} picks the epoch and restores its weights. With class
weights, the loss weights follow the same law as in the trees, and the power is searched. This
network is also the network of the two LTN arms.

The 1D-CNN reads the transformed vector of a flow as a sequence with one channel, as long as the
number of inputs. Two convolutions of 32 and 64 filters of width 3, with ReLU and a padding that
keeps the length, are followed by max pooling of size 2, flattening and dropout 0.2, then a
dense layer of 128 units with ReLU and dropout 0.2, and a softmax output. The weights start
from seeded Glorot uniform draws, as in the MLP. The class reuses the training loop of MLP-CE
unchanged, with the same optimiser, batch size, epoch cap, selection rule and class weights,
and its search covers the learning rate and the weight power over the ranges of MLP-CE with
weights. The 1D-CNN has more parameters than the common network: 506\,821 against 65\,029 on
NSL-KDD, nearly all in the dense layer after the pooling, and 351\,431 on CIC-IDS2017. It
therefore lies outside matched capacity. The order of the columns follows the transformer and
carries no notion of neighbourhood, so in our view the convolution tests an architecture, not
a model of the traffic.

The LSTM is an extension and is not implemented. It needs flows ordered in windows, which the
repository does not build.

\subsection{NeSy: LTN with the domain knowledge base}
\label{sec:impl-ltn}

LTN-sup trains the same network with a different loss. The class predicates are the outputs of
the softmax. For each class $c$ there is the axiom $\forall x\in C_c\; P_c(x)$, whose degree is
the $p$-mean-error aggregate of the output on the rows of that class in the batch
\citep{badreddine2022ltn}. The degrees of the classes are aggregated again with the same
exponent $p$, and the loss is one minus this satisfaction. Batches are not drawn uniformly. The
quota of a class is proportional to its count raised to $1-\mathrm{power}$, never below a
floor of rows per class, and a class smaller than its quota is drawn with replacement. The
power and the floor are searched. An epoch is as many batches as a pass over the training
rows would need, and it does not show every distinct row.

LTN-full adds one axiom per rule of the knowledge base, $\forall x\,(A_j(x)\rightarrow
P_{t_j}(x))$, over all rows of the batch. The antecedent $A_j$ is the product of its
conditions. A comparison is graded: its degree is the sigmoid of the distance to the threshold
over a width $s$ given in the file and multiplied by one factor that is searched. An equality is
crisp. The implication is the Reichenbach one, $1-a+a\,P$. The supervised and the domain
axioms are aggregated apart and combined with a weight $w$:
\begin{equation*}
\mathcal{L} = 1-\bigl((1-w)\,\mathrm{SatAgg}_{\mathrm{sup}}+w\,\mathrm{SatAgg}_{\mathrm{dom}}\bigr).
\end{equation*}
The antecedents are computed once on the raw training rows and indexed with each batch. The
crisp setting replaces every graded condition by a threshold and is used only in E4. At
inference both LTN arms run the network alone, so the knowledge is not evaluated.

\subsection{[EXT] Semantic-loss variant and/or post-hoc symbolic filter, as NeSy controls}
\label{sec:impl-controls}

The semantic-loss variant is not implemented, and no result of this thesis comes from it. The
post-hoc filter is. It is the post-check of the knowledge base on the raw row of a flow, after
any model has decided, and has two modes. The annotating mode leaves the prediction as it is and
finds the axioms that fire on each alert. The enforcing mode replaces the prediction by the
target of the axiom that fires, with the first axiom of the file winning a conflict, as in the
rule-only arm. Rules are read crisp, at the threshold 0.5 of C5. The module computes the
explanation measures of \cref{sec:explanation-quality} over the alerts: coverage, faithfulness
and conciseness. It runs inside a run on the validation or the test part, and on saved models
(\cref{sec:impl-harness}).

% ================================================================
% 9.4 Batch experimental harness
% ================================================================
\section{Batch experimental harness}
\label{sec:impl-harness}

\Cref{sec:harness} states what the harness must do. This section says what the code does.

\paragraph{Runner.}
One configuration and one seed make one result file, named after the run and the seed. The
runner loads the dataset, cuts the split, fits the transformer and the model, and scores the
validation part. It scores the test part only when asked, and the flag is false in every base,
tuning and stress configuration; the launch scripts of the final and stress runs set it on the
command line. For each scored part it computes the accuracy, the macro-F1, the MCC, the
precision, recall, F1 and support of each class and the confusion matrix. It adds the PR-AUC of
each class, the operator view of attack against Normal, the recall by fine label, the C5 report
of each axiom at the threshold 0.5 and, if asked, the post-check in one or both modes. Finished
seeds are skipped on resume, and a result file made with another configuration is an error. A
run can save the model and the preprocessing.

\paragraph{Result file.}
The file holds the configuration, the metrics of the parts, the timings of the stages, the
record of the training with the validation history, the epoch kept and the number of
parameters, the rows per class of each part and the sizes of the saved files. It also records
the commit, a flag for uncommitted changes, the versions of the packages and the platform. The
version of the knowledge base is in the record of the training and in the C5 report.

\paragraph{Stress runs.}
A generator writes the configurations of a stress experiment from a tuned configuration and
changes nothing else, so no hyperparameter is searched again. \Cref{tab:impl-stress} lists what
it varies and the levels it uses. A stress configuration is a normal configuration, and the
same runner executes it with the test part scored. The 1D-CNN has no stress configuration.

\begin{table}[t]
\centering
\footnotesize
\caption{What the generator of the stress runs varies, with the levels written in the code.}
\label{tab:impl-stress}
\setlength{\tabcolsep}{3pt}
\begin{tabular}{@{}p{1.2cm}p{4.2cm}p{7.3cm}@{}}
\toprule
 & Varied & Levels \\
\midrule
E3 & axioms removed & each axiom, then each group, when the base has more than one \\
E4 & LTN settings, one at a time & $p\in\{1,2,4,6\}$ other than the tuned one; for LTN-full crisp thresholds, width factor 0.5, 2 and 4 times the tuned one, $w\in\{0.1,0.25,0.5,0.75\}$ other than the tuned one \\
E7 & fraction of training and validation & NSL-KDD 0.5, 0.2, 0.1, 0.05; CIC-IDS2017 0.25, 0.1, 0.03, 0.01; CIC-DDoS2019 0.1, 0.01, 0.001; floor of 5 rows per class; for LTN-full also with the data axioms removed \\
E8 & classes made rarer in training and validation & R2L and U2R (NSL-KDD); R2L, U2R and Botnet (CIC-IDS2017); Normal (CIC-DDoS2019); kept 0.5, 0.25, 0.1; floor of 5 rows \\
E9 & fine labels held out & listed in the call, with their Attempted flows; all join the test part \\
Transfer & dataset and knowledge base & the CIC-IDS2017 base on \texttt{cic\_transfer}, validation 15\,\% of the source; knowledge base named in the call \\
E11 & thresholds and targets & shift of $-10$, $-5$, $+5$, $+10$, $+20$ percentile points of the training column; swap of the target of one axiom \\
E13 & copies of one axiom & 2, 4, 6, 8, thresholds moved on alternating sides by up to 10 percentile points \\
\bottomrule
\end{tabular}
\end{table}

\paragraph{Launchers.}
Shell scripts drive the runs, with seeds or trials as parallel processes and finished ones
skipped on restart. \texttt{tune\_final.sh} runs the search of one arm on one dataset, writes
the best configuration and calls the script of the final runs, which first checks the
configuration, the GPU and the row counts of the data, then runs the ten seeds with the test part
scored and the models saved, and archives results and models with a checksum.
\texttt{stress\_run.sh} runs a set of stress configurations. Before any run it reads each
dataset and reports its rows, and checks that the held-out labels exist in the data, that the
classes to be thinned exist, that each knowledge base loads with its edit applied and that its
columns are in the data. If a check fails, no run starts. The same script can stop after the
checks or save the models of the runs, which the transfer needs for the post-check.
\texttt{saved\_runs.sh} runs the tools on saved models for every arm of a dataset, and
\texttt{e13\_time.sh} times E13: it writes a configuration for the base and for each number of
copies, with five epochs and no early stop, runs them in one process over seeds 0 to 2 with the
seed in the outer loop, and prints the median time per epoch without the first epoch and its
ratio to the base.

\paragraph{Experiments on saved models.}
\texttt{saved\_eval.py} rebuilds the split of a final run from its configuration and seed,
applies the saved preprocessing, feeds the saved model and scores a part. The rule-only arm is
scored from its saved knowledge base on the raw rows. The tool prints the macro-F1 stored by the
run beside the new one. It scores the model alone and after the post-check in both modes, with
the knowledge base of the C5 report unless another is given. For each mode it records the
macro-F1, the predictions changed, the overrides that correct a wrong prediction, those that
break a correct one and the explanation measures; it also records, for each axiom, the rows
where it fires and its precision on the part, and the recall by fine label. The same tool
applies the perturbations of E10 to the raw part (\cref{sec:batch-protocol}). The noise uses
the seed of the run and the range of each column in the training part. The evasion finds, for
each axiom, its first comparison condition and, on the rows of its target class where it
fires, sets the column to $\tau-D\,|\tau|$ for a condition $>\tau$ and to $\tau+D\,|\tau|$ for a
condition $\le\tau$, with $|\tau|$ read as 1 when $\tau=0$; an axiom with equality conditions only
is left alone, and the tool counts the rows moved. The script runs noise at 0.1, 0.25 and 0.5 and
evasion at 0.1.

Launched by that script, \texttt{latency.py} times the saved model of seed 0 on the CPU, on rows of the test part already
preprocessed, with 20 warm-up calls and 200 timed calls for batches of 1, 64 and 1024. It
reports the percentiles 50, 95 and 99, the mean and the rows per second, and the size of the
model files. With the post-check it times, on the same batches, the network alone, the network
with the rules in annotate mode and the network with rules and override in enforce mode; the
rules read the raw columns of the same rows. The rule-only arm has no model file and is not
timed.

\paragraph{Analysis and scripts.}
The comparison module holds the tests of \cref{sec:statistical-plan}: the Wilcoxon test on
seeds, the Holm correction, the bootstrap interval of the median difference, and the Friedman
and Nemenyi tests across arms. Scripts build the tables from the result files, so no table of
this thesis is typed by hand. \texttt{stress\_curves.py} gives E7 and E8 against the rows
trained on, read from the result files and estimated for the final runs from the run that keeps
the largest share of each class. It reports the knowledge effect, LTN-full minus LTN-sup paired
by seed, and can read LTN-full from the runs with the data axioms removed.
\texttt{stress\_tables.py} gives E3, E4, E9, E11 and E13. For E9 it reports, over the held-out
labels, the share detected and the recall, and the share of their Attempted flows detected, which
are false alarms; the same measure on the final runs gives the labels found only in the test
part. \texttt{saved\_tables.py} gives the post-check, E10, E17 and E12 from the reports of the
tools on saved models. Differences are paired by seed, and none of these tables applies the Holm
correction.
% FIXME(open): no script builds the table of the transfer of E9 yet; summarize_runs.py reads its result files like any run. Add the table or name the command in Ch. 10.

% ================================================================
% 9.5 Real-time pipeline implementation
% ================================================================
\section{Real-time pipeline implementation}
\label{sec:realtime-pipeline}

The streaming pipeline is the work of the co-supervisor. The repository holds none of it: no
replay tool, extractor, broker, service, timer or dashboard. This section gives the design that
\cref{sec:streaming-pipeline} and \cref{sec:streaming-protocol} fix, and what the batch side
hands over. The rest is open, and no measurement of E14 or E15 exists.
% FIXME(open): the stack of the co-supervisor (Kafka and Spark Structured Streaming with the inference inside Spark, as reported orally) is in no file of the thesis and conflicts with the single-process pipeline of Ch. 7, Sec. 7.7.2, which is the mandatory part. Get the decision in writing and align this section, Ch. 7 and Ch. 8.

\subsection{Traffic source: pcap replay with rate control}
\label{sec:impl-source}

The source is the replay of a captured dataset at its recorded rate and at multiples of it. A
live capture is the alternative of the design. The replay clock is the clock of the capture,
on which the attack start of E15 is measured.
% FIXME(open): replay tool, list of speeds, which captures (the pcap files of CIC-IDS2017 and of CIC-DDoS2019 are not in the data of this repository) and whether live capture is used.

\subsection{Flow assembly and feature extraction}
\label{sec:impl-flows}

Flow assembly closes the flows, so that the rest of the pipeline sees whole flows. The model of
the pipeline is the one trained on the corrected CIC-IDS2017, with the classes of
\cref{tab:impl-classes} and an alert on the DDoS class. Since the batch code computes no window
feature, the vector is the per-flow columns and the protocol, scaled and coded by the saved
transformer. The message carries the raw record beside the vector, because the symbolic check
reads raw columns.
% FIXME(open): confirm the extractor, the model and the class served (Ch. 7, Sec. 7.7.2). The corrected CICFlowMeter of Engelen et al. is not packaged here, and its exporter timeout is the first part of the detection delay.

\subsection{Message broker and consumer topology}
\label{sec:impl-broker}

A broker sits between producer and consumer and absorbs bursts. Its queue is bounded. Under
backpressure the producer is slowed in one condition and flows are dropped in another, and the
share of dropped flows is reported with the throughput. The mandatory topology is one inference
process. A distributed engine is an extension.
% FIXME(open): choice of the broker, size of the queue, the two backpressure conditions, number of consumers, and whether DDL is used (if so, amend the scope paragraph of Ch. 1, Sec. 1.6).

\subsection{Inference service and optional symbolic post-check}
\label{sec:impl-service}

The service loads an exported model and calls it on micro-batches. A final run saves, per seed,
the network as a \texttt{.keras} file, the 1D-CNN included, a forest as a \texttt{joblib} file
or a booster as JSON with the rounds kept, and the fitted transformer with the feature and class
names. The latency tool already loads these three kinds. The post-check is a function of the raw
record and the scores, and the service can call it after the model; its cost on single rows and
small batches is the one E12 measures. The rule-only arm needs the knowledge base file and no
model. The service itself does not exist. FR9 is checked by the share of identical decisions in
batch and in streaming.
% FIXME(open): the service, its input format, the micro-batch size and the check of FR9 are not written.

\subsection{Instrumentation: timestamps, metric collection, dashboard}
\label{sec:impl-instrumentation}

Four stamps are taken for each flow: when it is closed, when it reaches the broker, when it is
decided and when it is alerted. Their differences give the exporter timeout, the transport, the
inference and the check. The detection delay of E15 is the time of the first alert on a DDoS
flow minus the start of the first DDoS flow of the replay.
% FIXME(open): none of the stamps, the collection of metrics or the dashboard is implemented. State which measures belong to E14 (throughput, latency distribution and share of dropped flows under each speed) and which to E15 (detection delay); Ch. 8 names only E15, and Ch. 10 asks for the division.

% ================================================================
% 9.6 Hyperparameter search protocol
% ================================================================
\section{Hyperparameter search protocol}
\label{sec:impl-tuning}
A tuning file names a base configuration, the number of trials, a search seed, the tuning
seeds, the settings fixed for all trials, the settings fixed for the trials only, and the
space, in which each setting is a list of values, a uniform range, a log-uniform range or an
integer range (\cref{app:hyperparameters}). The trials are the 20 points of a Latin hypercube
over these settings, drawn with search seed 0. The grid depends on the space and the seed only,
so the same arm gets the same 20 settings on every dataset. Trial $k$ of an arm is the same
setting on the three datasets, and when it wins on two of them the two tuned configurations
coincide. This happened for Random Forest on the two CIC datasets (trial 9), for MLP-CE with
weights on the same two (trial 5) and for LTN-full on NSL-KDD and CIC-DDoS2019 (trial 6), while
XGBoost won with trial 10 on one CIC dataset and trial 4 on the other. A likeness of the
settings across datasets is therefore not independent evidence.

A trial is the base configuration with the settings drawn, run on the tuning seeds 100, 101 and
102, which are not among the final seeds. Its score is the validation macro-F1 of the model
kept by the stopping rule, averaged over the three seeds. The best trial has the highest mean,
and a tie goes to the lowest index. A loss that becomes non-finite fails the trial, which is
recorded as a result. The test part is never scored: the routine refuses a base configuration
that scores it, and it drops the knowledge base of the C5 report. Finished runs are reused, so
an interrupted search resumes, and a result file made with another configuration stops it. The
trials can be split over parallel workers, worker $i$ taking the trials with index $i$ modulo
the number of workers, and a last step writes the summary and the best configuration, a normal
configuration with the chosen settings and the ten final seeds.

The fixed settings of the four neural arms are those of \cref{sec:impl-mlp}: hidden layers of
256 and 128 units, dropout 0.2, batch size 256, Adam and at most 100 epochs, with the stopping
rule of the base file. The 1D-CNN fixes its own network and the same batch size and epoch cap.
In the search on CIC-IDS2017 the training part holds a quarter of the Normal rows, drawn at
random, as a setting of the trials only. The validation part stays whole and the final runs use
all training rows, so a validation score of the search is not comparable with one of a final
run.

Three limits follow from the results of the searches.
\begin{itemize}
\item The epoch cap hardly binds. Of the 60 runs of an arm, 0 to 2 end at 100 epochs: two for
  LTN-sup and one for MLP-CE with weights on NSL-KDD, one for MLP-CE without weights on
  CIC-DDoS2019, none elsewhere. The best trial of every arm ended below the cap on all three
  seeds, and none of the 10 final runs of a neural arm on the CIC datasets ends at the cap.
\item Several best trials lie on the edge of the space: the number of trees and the leaf size
  of Random Forest on both CIC datasets, the learning rate and the column sampling of XGBoost on
  CIC-DDoS2019, and for LTN-full the quota floor, the exponent $p$ and the domain weight on
  CIC-DDoS2019 and the learning rate, the width scale and the domain weight on CIC-IDS2017.
  The best setting may lie outside the space. The search was not repeated, since the budget is
  fixed, and Appendix~\ref{app:hyperparameters} declares it as a limit.
\item The budget is equal in trials and not in dimensions. MLP-CE without weights has one
  setting to search, MLP-CE with weights and the 1D-CNN two, LTN-sup four, Random Forest five,
  XGBoost and LTN-full six (Appendix~\ref{app:hyperparameters}).
\end{itemize}
% FIXME(open): number of NSL-KDD final runs that end at the epoch cap, from epoch_cap_report.py on the final files; Ch. 10 asks for it too.
% FIXME(open): the runs at the epoch cap and the edges of the space for the search of the 1D-CNN, once its tuning files are back.

% ================================================================
% 9.7 Reproducibility package
% ================================================================
\section{Reproducibility package (repository structure, seeds, environment file, run scripts)}
\label{sec:impl-repro}

The package is the repository \texttt{fil2405/nesy-nids-bench}. The library is in
\texttt{src/nesybench}, the configurations of arms, tuning and stress runs in \texttt{configs},
the knowledge bases in \texttt{kb}, the analysis scripts in \texttt{analysis} and the launch
scripts in \texttt{scripts}. The downloaded files, the cleaned tables and the result files go
to \texttt{data} and \texttt{results}, which are not under version control. The directory
\texttt{tests} has 37 files with 395 test functions.

A run needs a configuration and a seed. The seeds of the final runs are 0 to 9 and those of the
tuning are 100 to 102. A seed sets Python, NumPy and TensorFlow, and it switches on the
deterministic operations of TensorFlow. The package also sets \texttt{TF\_ENABLE\_ONEDNN\_OPTS}
to 0 when the environment leaves it unset, so a run that starts from the launch scripts has it
at 0 unless the shell exports another value. The test cut of CIC-IDS2017 has its own seed, 0,
and the validation cut follows the seed of the run (\cref{sec:impl-splits}).

A second run of the same identifier and seed is refused unless the caller asks to overwrite.
All tables of the thesis are produced from the result files by the scripts of \texttt{analysis},
so a number in a table can be traced to the runs behind it. The environment file pins the
versions of the laptop; those of the server are known from the result files
(\cref{tab:impl-environments}).

Two limits apply to the runs of the thesis. On the server the code was not a repository, so
the tuning files, the final runs and the stress runs record the commit as \texttt{unknown} and
leave the flag of the modified tree empty. A series is identified by its date and by the
configuration stored in each file. The second limit is the environment: determinism holds
within one, as \cref{sec:impl-environment} says.
% FIXME(open): record the commit of each series before any figure is called citable (open point of Ch. 8, Sec. 8.1.5), and check that all finals used the same TF_ENABLE_ONEDNN_OPTS.
% FIXME(open): the environment file of the server (TensorFlow 2.22.0-rc0, Python 3.14) is not in the repository; decide whether to add a second lock file and where the result files are archived (Appendix E).
% FIXME(open): the streaming pipeline has no code, no configuration and no test in the repository; add its directory and its launch script when the co-supervisor delivers them.

% ================================================================
% 9.8 Summary
% ================================================================
\section{Summary}
\label{sec:impl-summary}

The batch benchmark is complete for the arms of the plan that exist: the rule-only classifier,
Random Forest, XGBoost, MLP-CE with and without weights, the 1D-CNN, LTN-sup and LTN-full, with
the post-check in both modes. Around them sit the runner, the search, the stress generator with
the schema of the transfer, the launchers with their checks, the tools on saved models, the
latency tool, the perturbations and the scripts that build the tables. The LSTM and the
semantic-loss variant are not implemented, and the thesis does not report them. The streaming
pipeline is not part of the repository, and \cref{sec:realtime-pipeline} gives its design only.
